% Chương 1: Giới thiệu
\chapter{GIỚI THIỆU}

\section{Giới thiệu đề tài}

Đề tài ``Nền tảng kết nối cộng đồng nhiếp ảnh phim và các Film Lab (MiniLab) tích hợp AI'' hướng đến kết nối người chụp ảnh phim, cơ sở xử lý phim, chuyên gia nhiếp ảnh và đối tác giao nhận trên cùng một hệ thống. Nền tảng hỗ trợ các dịch vụ tráng phim, quét phim (scan), in ảnh, đồng thời tạo không gian lưu trữ ảnh số, trao đổi thiết bị và chia sẻ kiến thức nhiếp ảnh.

\section{Mục đích của đề tài}

Nền tảng giúp người dùng thuận tiện tìm kiếm Film Lab, đặt dịch vụ và theo dõi quá trình xử lý phim; giúp Film Lab quản lý đơn hàng, khách hàng và bàn giao ảnh số tập trung. Trong phạm vi đồ án, nhóm tập trung thiết kế cơ sở dữ liệu PostgreSQL để tổ chức, lưu trữ và quản lý dữ liệu phục vụ các nghiệp vụ này, đồng thời hỗ trợ các chức năng cộng đồng và tích hợp AI.

\section{Các chức năng chính}

\begin{itemize}
    \item Tìm kiếm, so sánh và đánh giá Film Lab theo vị trí, dịch vụ, giá và phản hồi của khách hàng.
    \item Đặt dịch vụ tráng, scan, in ảnh; yêu cầu giao nhận, thanh toán trực tuyến và theo dõi tiến độ đơn hàng.
    \item Quản lý thông tin Film Lab, gói dịch vụ, khách hàng, quy trình xử lý phim và doanh thu.
    \item Bàn giao, tải xuống và tổ chức ảnh scan trong kho ảnh phim số cá nhân, kèm thông tin về phim và thiết bị chụp.
    \item Mua bán, trao đổi máy ảnh, ống kính, phim và phụ kiện; chia sẻ bài viết, thảo luận và tham gia sự kiện nhiếp ảnh.
    \item Tích hợp AI để gợi ý Film Lab và loại phim, hỗ trợ tìm kiếm, phân tích chất lượng ảnh scan và giải đáp kiến thức nhiếp ảnh.
    \item Quản trị người dùng, phân quyền, phê duyệt Film Lab, kiểm duyệt nội dung và giải quyết khiếu nại.
\end{itemize}
